\chapter{Gen Z Đi Làm Sớm}
\markboth{Gen Z Đi Làm Sớm}{Gen Z Đi Làm Sớm}

\begin{truyen}{Ra Ngoài Đời Không Có Điểm Cộng Cảm Xúc}{Early Work Reality Check}
\chuhoa{C}{ó bạn mới đi làm thêm được hai tuần đã than với thầy: ``Em thấy môi trường làm việc toxic quá. Sếp giao việc gấp mà còn khó tính. Em tính nghỉ, vì em không muốn đánh đổi sức khỏe tinh thần.''}

\teacherpair{Thầy hỏi ngay: ``Sếp bắt em làm sai đạo đức, xúc phạm em, hay chỉ đơn giản là yêu cầu đúng hạn và góp ý thẳng?''}{The teacher asks: ``Is your boss making you do immoral things, insulting you, or simply demanding deadlines and giving blunt feedback?''}

Bạn ấy ngập ngừng: ``Dạ... chủ yếu là góp ý thẳng và bắt đúng giờ.''

\teacherpair{Vậy thì có khi em không đang gặp môi trường độc hại. Em chỉ đang gặp thế giới không có điểm cộng vì em dễ thương. Ở trường, nhiều khi thầy cô còn nhắc em mấy lần. Ra ngoài, người ta trả tiền để em làm được việc, không phải để ôm cảm xúc của em mỗi ngày.}{``Then maybe you are not meeting a toxic environment. You are meeting a world that gives no bonus points for being likable. At school, teachers may remind you many times. Outside, people pay you to do the work, not to hold your emotions every day.''}

Thầy dịu giọng lại một chút: ``Dĩ nhiên, nếu bị xúc phạm, bóc lột, hay ép làm điều sai thì phải biết bảo vệ mình. Trưởng thành không có nghĩa là chịu đựng vô điều kiện. Nhưng cũng đừng vì chưa quen bị góp ý mà gắn nhãn độc hại cho mọi đòi hỏi bình thường. Nếu không phân biệt được hai thứ đó, em sẽ vừa thiệt thân vừa khó lớn.''

Bạn ấy im rất lâu. Kiểu im của người bắt đầu hiểu rằng nhiều lý sự trong lớp mang ra xã hội sẽ vỡ rất nhanh.

\textit{(A student calls ordinary workplace pressure toxic. The teacher distinguishes abuse from professional expectation.)}
\end{truyen}

\begin{baihoc}
Đi làm sớm là tốt nếu nó giúp mình lớn lên. Còn nếu chỉ mang nguyên lý sự học đường ra xã hội, mình sẽ vỡ mộng rất sớm.
\textit{(Working early is good if it helps you grow. If you carry classroom excuses straight into society, disillusion comes quickly.)}
\end{baihoc}

\begin{ghinhoanh}
\textbf{Short English to remember:}

Not every pressure is abuse.

Work expects results, not excuses.
\end{ghinhoanh}

\begin{tuhocnhanh}
1. Em có đang gọi mọi góp ý khó nghe là độc hại không? \textit{(Are you calling every uncomfortable feedback toxic?)}

2. Nếu bước vào môi trường làm việc ngày mai, em cần bỏ lại lý sự nào ở cổng trường? \textit{(If you entered a workplace tomorrow, which excuse should you leave at the school gate?)}
\end{tuhocnhanh}

\ngancach